\SourceLeaf{153}sionen bestehen, sonst aber die auf die einzelnen Pensionen gleichmässig zu vertheilende Gesamtsumme von 150 Rthl. Dasselbe gilt von jeder folgenden Erhöhung.
\end{quote}
\noindent Auf diese Weise hätte man gar nicht einmal nöthig gehabt, den \textit{Anfang} einer Erhöhung von der Bedingung einer nicht über 15 hinausgehenden Pensionenzahl abhängig zu machen. Und in dieser Beziehung hätte diese Bestimmungsart sich sogar als vortheilhafter für die Witwen gezeigt, weil bei einer andauernd bestehenden kleinen Überzahl der Verlust, welcher aus der Verpflichtung entsteht, die bisherige Erhöhungssumme mit mehrern theilen zu müssen, bald durch eine neue Erhöhung compensirt oder mehr als compensirt würde, während in einem solchen Fall eine neue Erhöhung nach der P.’schen Normirung und in der zweiten Interpretation gar nicht zulässig ist.

Ich möchte übrigens glauben, dass die obige abgeänderte Fassung, wenn damals jemand daran gedacht hätte, sie in Vorschlag zu bringen, auch P. hätte zufrieden stellen müssen. Denn entweder war seine Voraussetzung, 15 sei die höchste Pensionenzahl, die vorkommen könne, richtig, oder sie war unrichtig: im ersten Fall war die Abänderung ganz wirkungslos, mithin gleichgültig, im zweiten aber nothwendig.

Die Folge einer solchen Normirung wäre gewesen, dass man sich gewöhnt haben würde, die Pension wie aus zwei Theilen zusammengesetzt zu betrachten, einem unveränderlichen Theile und einer Zulage, die von Zeit zu Zeit mit dem Kapitalvermögen wachsen, möglicherweise aber auch dabei wieder etwas zurückgehen könne, letzteres aber dann nach einer wirklich unwandelbaren einfachen Regel, wonach jede Witwe leicht selbst die Controle führen könnte. Auf den Unterschied zwischen einem solchen gesetzlichen Zurückgehen, und dem, nach der Quittungs-Clausel in Folge des Unvermögens der Kasse eintretenden habe ich schon oben S. [131] aufmerksam gemacht.

Drittens scheint es wohl der Mühe werth zu sein, die Ursachen anzugeben, welchen man das rasche und ununterbrochene Steigen der Prosperität der Witwenkasse während eines Zeitraums von mehr als vierzig Jahren zuzuschreiben hat. Ich finde, dass diess Steigen ganz vorzüglich begünstigt ist durch das Zusammenwirken von zwei Umständen, auf deren einen 1794 gar nicht gerechnet war, auf den andern aber wenigstens nicht gerechnet werden durfte.

Die erste Ursache ist das Steigen des Zinsfusses, welches schon wenige Jahre nach jener Epoche anhob. Man hatte, wie ich nachgewiesen habe, damals mit Sicherheit nur auf 3 Proc. rechnen zu dürfen geglaubt. Aber schon 1799 bemerkte das Curatorium in dem schon oben angeführten Rescript vom 23. Mai, dass Gelegenheit zu sicherer Unterbringung zu 4 Proc. gar nicht selten sei, und bot sogar die eigne Mitwirkung dazu an. Etwas später aber erhob sich der Zinsfuss allgemein auf 5 Proc., und beharrte für den grössten Theil der Kapitalien der Witwenkasse während einer langen Reihe von Jahren auf dieser Höhe.

Die zweite Ursache ist der Umstand, dass die Zahl der Pensionen während jenes langen Zeitraums \textit{unter} dem zu erwartenden Mittelwerthe (15) geblieben ist, ja man muss sagen \textit{beträchtlich} unter dem zu erwartenden Mittelwerthe, wenn man dafür nach dem plausiblern Anschlage S. [145] die Zahl 17 annimmt (nemlich \(\tfrac{4}{7}\times26\) Witwen mit Zusatz von \(\tfrac{1}{6}\) wegen der Waisen). Dass man diese Erscheinung gar nicht wie etwas sehr ausserordentliches zu betrachten habe, ist schon oben S. [144 Z. 7] erwähnt: allein eben so wenig wäre es etwas ausserordentliches gewesen, wenn gerade das Gegentheil eingetreten, und z. B. schon 20—25 Jahre nach jener Epoche von 1794 die Mittelzahl auf der andern Seite bedeutend und andauernd überschritten wäre. In einem solchen Falle würde an die Stelle des raschen Steigens des Vermögens der Kasse ein sehr langsames getreten sein, ja vielleicht ein Stillstand oder sogar die Nothwendigkeit zu der Quittungsklausel die Zuflucht zu nehmen, falls sich zugleich ein tiefes Herabgehen des Zinsfusses dazu gesellt hätte.

Höhere Witwenzahl ist nun bereits seit einer Anzahl von Jahren eingetreten S. [126 Z. 23] und
